\scanpage{154}
toute la composante connexe du dit point). Soit $P$ son complémentaire. Il est \inserted{un rev. étale de $Y$,} stable sous $\unsure{\mathfrak S_n|\mathfrak G_n|G_n}$, qui y opère sans inertie. De plus, en tout p\textsuperscript{t} $y\in Y$, l'\deleted{\unsure{ensemble des|étale}}\inserted{fibre} $P$ \unsure{sur|en} $y$ a exactement $n!$ points géométriques --- car \deleted{\unsure{elle est}} celle $E$ de $\deleted{\unsure{X^n}}X$ est isomorphe à $I_n$, donc celle de $X^n$ est \deleted{isomorphe} \inserted{isomorphe} comme ens. à opérateurs à $E^n$ \deleted{\unsure{I_n^n}}, et on \unsure{est|a} à enlever les pts qui sont fixes sous des permutations non triviales de $I_n$ : \unsure{il reste|ce sont} les suites d'éléms \uline{tous} \uline{distincts} de $E$, \inserted{il y en a $n!$.}

En vertu de prop 2.5., $P$ est \inserted{donc} un revêtement étale galoisien de $Y$ de groupe $\mathfrak S_n$.

D'autre part, \deleted{on sait} la \deleted{première} \inserted{dernière} projection $\mathrm{pr}_{\unsure{12|n}}:X^n\longrightarrow X$ \deleted{\unsure{est inv|étale}} est invariante sous le stabilisateur $H=\mathfrak S_{n-1}$ de $\{n\}$, \deleted{d'où un hom} il en est à fortiori du morphisme induit sur $P$ : $P\longrightarrow X$, d'où un \deleted{isomorphisme} $Y$-\deleted{hom}morphisme canonique

\[
\varphi:P/H\longrightarrow X
\]

\margin{?\\ \deleted{\unsure{?}}}

Comme $P/H$ est étale sur \deleted{\unsure{$Y$}} (\uline{vérif}) \inserted{ainsi que $X$ \unsure{c'est|est} un morphisme étale} pour montrer que $\varphi$ est un isom, il suffit de le prouver

% archive sheet number: 153

